% ============================================================================
% Fichier  : partie1/P01_C01_controler_auditer_si_public.tex
% Titre    : Contrôler et Auditer un SI Public
% Auteur   : MINKA MI NGUIDJOÏ Thierry Emmanuel
% Version  : 0.1 -- novembre 2026
% ============================================================================
% Sources : INF3226 C01 (resserré) ; ISSAI 100 (via le manuel IDI) ;
% GUID 5100 ; décret n° 2013/287, art. 2. Seuls A0 et A1 (HAL) sont cités
% comme travaux de l'auteur.
% ============================================================================

\chapter{Contrôler et Auditer un SI Public}
\label{chap:audit-public}

\epigraph{%
    « Auditer, c'est donner à un tiers une assurance sur ce qu'il ne peut
    pas vérifier lui-même. »%
}{Principe du cours}

\niveauChapitre{Fondamental}{Prologue. Aucune connaissance informatique
préalable.}

\begin{boxobjectifs}
\begin{itemize}
    \item Distinguer contrôle, audit, inspection et vérification, et situer
          chacun dans l'administration financière camerounaise.
    \item Expliquer la différence entre assurance raisonnable et assurance
          limitée.
    \item Reconnaître les trois types d'audit du secteur public (financier,
          de conformité, de performance) et ce que chacun demande au
          système d'information.
    \item Situer le contrôle des systèmes d'information parmi les cinq
          contrôles du CONSUPE.
    \item Cartographier les systèmes d'information des finances publiques et
          leurs points de contrôle.
    \item Expliquer ce que l'informatique change à la preuve.
\end{itemize}
\end{boxobjectifs}

\bigskip

% ============================================================================
\section*{Pourquoi ce chapitre}
% ============================================================================

Les auditeurs du PSSFP connaissent déjà le contrôle des finances
publiques. Ce chapitre ne le leur réapprend pas~: il fixe un vocabulaire
commun et montre ce qui change quand l'objet du contrôle devient un système
d'information. La question directrice est simple~: \emph{que faut-il
regarder, et dans quel ordre, quand la dépense passe par une
application~?}

\separateur

% ============================================================================
\section{Contrôler, auditer, inspecter, vérifier}
\label{sec:c01-notions}
% ============================================================================

Les mots du contrôle sont employés indifféremment dans les services. Le
cours les distingue, parce que chacun engage une démarche et une
responsabilité différentes.

\begin{boxdef}
\termecle{Contrôle}\index{contrôle}~: dispositif ou action, intégré au
fonctionnement de l'entité, qui vise à empêcher ou à détecter une anomalie
(un visa, un rapprochement, un blocage automatique). Il est exercé par les
services eux-mêmes.

\medskip
\termecle{Audit}\index{audit}~: examen méthodique, indépendant et
documenté, qui compare une situation à des critères connus à l'avance et
qui conduit à une opinion motivée, communiquée à un tiers.

\medskip
\termecle{Inspection}\index{inspection}~: vérification ponctuelle de la
régularité d'opérations ou du fonctionnement d'un service, souvent
déclenchée par un événement, au nom de l'autorité hiérarchique.

\medskip
\termecle{Vérification}\index{vérification}~: dans la langue du décret
n°~2013/287, l'activité des vérificateurs du Contrôle Supérieur de l'État,
conduite sur pièces et sur place~; elle recouvre l'audit externe du
secteur public.
\end{boxdef}

La différence tient en deux questions. \textbf{Qui agit~?} Le contrôle est
exercé par l'entité~; l'audit et la vérification par quelqu'un qui n'a pas
participé aux opérations. \textbf{Pour qui~?} Le contrôle sert l'entité~;
l'audit sert un tiers (le Président de la République pour le CONSUPE, le
juge des comptes, le parlement, un conseil d'administration) qui ne peut
pas vérifier lui-même.

\subsection{Les éléments d'un audit}

Tout audit, quel que soit son objet, manipule les mêmes cinq éléments. Le
cours les emploie dans ce sens précis.

\begin{tableaucours}{L{3.0cm}L{9.2cm}}{%
    \tableauHeader{Élément} &
    \tableauHeader{Définition et exemple (dossier MINÉQUIR)}
}
Critère & Exigence de référence, connue avant le test~: «~un mandat ne peut
pas être payé deux fois~» \\
Preuve & Information vérifiable que l'auditeur obtient et conserve~:
l'extraction des 9\,611 mandats, son empreinte \\
Constatation & Comparaison de la preuve au critère~: «~le mandat 2025-09342
reprend la facture du mandat 2025-08817~» \\
Conclusion & Résultat de l'ensemble des constatations sur un objectif~:
«~le contrôle des doublons n'existe pas dans SIDÉPENSE~» \\
Opinion & Forme que prend la conclusion dans le rapport, avec son niveau
d'assurance \\
\end{tableaucours}
\vspace{-0.8em}
\begin{center}{\small\itshape Tableau~1.1. Les cinq éléments d'un
audit}\end{center}

\begin{boxalert}
\textbf{Violation typique.} Choisir le critère après avoir vu les faits~:
«~les droits de M.~MBARGA étaient excessifs~», sans dire quelle règle les
limitait.

\smallskip
\textbf{Règle~:} le critère est écrit avant le test, et cité dans le
constat. Un constat sans critère est une opinion personnelle.
\end{boxalert}

\separateur

% ============================================================================
\section{Assurance raisonnable, assurance limitée}
\label{sec:c01-assurance}
% ============================================================================

Aucun auditeur n'examine tout. Il donne donc une assurance, jamais une
certitude.

\begin{boxdef}
\termecle{Assurance raisonnable}\index{assurance raisonnable}~: niveau de
confiance élevé, mais non absolu, obtenu par des travaux étendus. Elle
s'exprime sous une forme positive~: «~les contrôles examinés fonctionnent
efficacement~».

\medskip
\termecle{Assurance limitée}\index{assurance limitée}~: niveau de
confiance plus faible, obtenu par des travaux réduits. Elle s'exprime sous
une forme négative~: «~nous n'avons rien relevé qui laisse penser que…~».
\end{boxdef}

Le choix se fait au moment du mandat, pas à la fin de la mission. Une
inspection de deux jours ne peut donner qu'une assurance limitée~; si son
rapport conclut sous une forme positive, il promet plus qu'il n'a vérifié.
C'est exactement le défaut du rapport de janvier dans le prologue
(Chapitre~\ref{chap:prologue})~: «~régularité satisfaisante~» sur une
population incomplète.

\separateur

% ============================================================================
\section{Trois types d'audit dans le secteur public}
\label{sec:c01-types}
% ============================================================================

Les normes de l'INTOSAI distinguent trois types d'audit du secteur
public~\cite{idi2021conformite}. Chacun interroge le système d'information
d'une manière différente.

\begin{tableaucours}{L{2.6cm}L{4.6cm}L{5.0cm}}{%
    \tableauHeader{Type} &
    \tableauHeader{Question} &
    \tableauHeader{Ce qu'il demande au SI}
}
Audit financier & Les états financiers et les comptes sont-ils fiables~? &
Les données produites par l'application sont-elles exactes et
exhaustives~? \\
Audit de conformité & Les opérations respectent-elles les lois, les
règlements et les règles internes~? & L'application impose-t-elle les
règles (séparation des tâches, plafonds, visas)~? \\
Audit de performance & Les ressources sont-elles employées avec économie,
efficience et efficacité~? & Le système atteint-il ses objectifs, à quel
coût, et ses données permettent-elles de le savoir~? \\
\end{tableaucours}
\vspace{-0.8em}
\begin{center}{\small\itshape Tableau~1.2. Les trois types d'audit et le
système d'information}\end{center}

L'audit des systèmes d'information n'est pas un quatrième type à côté des
trois autres. Il les sert tous les trois~: on audite un SI pour savoir si
l'on peut se fier à ses données (financier), s'il fait respecter les règles
(conformité), ou s'il rend le service attendu (performance). Les lignes
directrices GUID~5100 de l'INTOSAI le disent ainsi~: un audit des SI peut
être une mission autonome, de performance ou de conformité, ou une partie
d'une mission plus large, financière, de conformité ou de
performance~\cite{intosai2019guid5100}.

\subsection{Les cinq contrôles du CONSUPE}

Le décret n°~2013/287 range le contrôle des systèmes d'information parmi
les cinq contrôles qu'effectuent les services du Contrôle supérieur de
l'État~\cite{decret2013287consupe}.

\begin{boxjuris}
\textbf{\loicam{Décret n°~2013/287}, article~2, alinéa~1}~: \textit{« Les
services du Contrôle supérieur de l'Etat constituent l'institution
supérieure de contrôle des finances publiques (ISC) du Cameroun. Ils sont
chargés de l'audit externe. »}

\medskip
\textbf{Article~2, alinéa~3}~: \textit{« Dans l'exercice de leurs
attributions, les services du Contrôle supérieur de l'Etat effectuent~:
un contrôle de conformité et de régularité~; un contrôle financier~; un
contrôle de performance~; un contrôle de l'environnement~; un contrôle des
systèmes d'information. »}

\medskip
\textbf{Portée pour l'audit.} Le décret reprend les trois types de
l'INTOSAI (conformité, financier, performance) et leur ajoute deux
contrôles, dont celui des systèmes d'information. Le vérificateur peut
donc conduire une mission dont l'objet principal est un SI, ou intégrer un
volet SI à une mission de régularité ou de performance.
\end{boxjuris}

\separateur

% ============================================================================
\section{Les systèmes d'information des finances publiques}
\label{sec:c01-cartographie}
% ============================================================================

\subsection{Une chaîne de systèmes}

Les finances publiques reposent sur une chaîne d'applications, de la
préparation du budget à la reddition des comptes. Le cours ne décrit pas
les applications réelles de l'État camerounais, qui évoluent et dont le
détail n'est pas public~: il raisonne sur des fonctions, que l'auditeur
retrouve sous des noms différents d'une administration à l'autre.

\begin{tableaucours}{L{2.6cm}L{4.8cm}L{4.8cm}}{%
    \tableauHeader{Fonction} &
    \tableauHeader{Ce que fait le système} &
    \tableauHeader{Risque typique}
}
Budget & Prépare, vote, notifie et modifie les crédits & Crédits modifiés
sans acte, virements non autorisés \\
Dépense & Engage, liquide, ordonnance, transmet au comptable & Double
paiement, service fait fictif, fractionnement \\
Recette & Liquide, émet et recouvre les titres de recettes & Titres
annulés sans motif, recouvrements non reversés \\
Paie & Calcule et verse la solde & Agents fictifs, éléments de solde
indus \\
Trésorerie & Exécute les paiements, tient les comptes de disponibilités
& Paiement sur un compte bancaire modifié \\
Comptabilité & Produit les comptes et les états financiers & États qui ne
correspondent pas aux opérations \\
\end{tableaucours}
\vspace{-0.8em}
\begin{center}{\small\itshape Tableau~1.3. Les fonctions des systèmes
financiers publics et leurs risques}\end{center}

\subsection{Les points de contrôle de la chaîne}

Dans une chaîne de systèmes, les anomalies naissent souvent aux
\textbf{jointures}~: là où une application transmet des données à une
autre (de la dépense vers la trésorerie, de la paie vers la comptabilité),
et là où une donnée de référence est modifiée (un RIB, un barème, un
plafond). L'auditeur y place ses premiers tests.

\begin{boxscenario}{Dossier MINÉQUIR. Trois jointures}
\textbf{Constatation.} Dans SIDÉPENSE, trois données conditionnent tous
les paiements~: le fichier des fournisseurs (et leurs RIB), la table des
profils d'utilisateurs, et le fichier de transmission des mandats au poste
comptable.

\medskip
\textbf{Analyse} (confiance élevée). Le double paiement du prologue a
utilisé deux de ces trois jointures~: un RIB modifié et un profil resté
actif. Une mission qui teste ces trois points en premier couvre une grande
part du risque de la chaîne, pour un effort limité.
\end{boxscenario}

\separateur

% ============================================================================
\section{Ce que l'informatique change à la preuve}
\label{sec:c01-preuve}
% ============================================================================

Le contrôle sur pièces repose sur une évidence~: la pièce existe, elle est
signée, elle est classée. Dans une administration informatisée, chacune de
ces évidences change de forme.

\begin{tableaucours}{L{3.2cm}L{4.4cm}L{4.6cm}}{%
    \tableauHeader{Sur papier} &
    \tableauHeader{Dans le système} &
    \tableauHeader{Ce que l'auditeur doit vérifier}
}
La pièce & Un enregistrement dans une base & Qu'il n'a pas été modifié
après coup \\
La signature & Un droit d'accès utilisé & Que le titulaire du compte est
bien celui qui a agi \\
Le visa & Un contrôle automatique & Que le paramétrage n'a pas été changé
\\
Le classeur & Une extraction & Qu'elle est complète, et qu'il l'a obtenue
lui-même \\
L'archive & Une sauvegarde & Qu'elle peut être restaurée \\
\end{tableaucours}
\vspace{-0.8em}
\begin{center}{\small\itshape Tableau~1.4. Ce que l'informatique change à
la preuve}\end{center}

Deux conséquences en découlent. D'abord, la \textbf{population} d'un test
est dans le système, pas dans les classeurs. Ensuite, une preuve
électronique ne vaut que si l'auditeur peut en établir l'intégrité~: une
extraction remise sur une clé, sans empreinte ni rapprochement, est une
déclaration de l'audité. Le Chapitre~\ref{chap:preuve} donne la manière de
faire.

\begin{boxinfo}
\textbf{Une réponse de plus.} Parce que la preuve électronique peut manquer
ou être incertaine, le cours admet trois réponses pour chaque contrôle
testé~: \OUI{}, \NON{} et \NEK{} (connaissance insuffisante). Elles
viennent d'un programme de recherche sur l'audit, l'auditologie, qui
cherche à mesurer ce que l'auditeur ne sait
pas~\cite{minka2026auditologie,minka2026nek}. Le cours n'en garde que
l'essentiel~: dire ce que l'on ne sait pas, plutôt que de le cacher sous un
\OUI{} ou un \NON{}.
\end{boxinfo}

\begin{boxmetier}{Vérificateur}
Avant la mission, demandez la liste des applications qui produisent les
opérations de votre périmètre, avec, pour chacune, le service
propriétaire et l'interlocuteur technique. Cette liste est la première
pièce de votre dossier permanent.
\end{boxmetier}

\begin{boxmetier}{Juge des comptes}
Un compte produit par une application n'est fiable que si les contrôles
de l'application le sont. Lorsque vous relevez une anomalie dans un
compte, demandez quel contrôle de l'application aurait dû l'empêcher~:
la réponse oriente l'instruction.
\end{boxmetier}

\begin{boxmetier}{Auditeur interne}
Votre rapport n'est pas lu que par votre direction~: l'article~5 du
décret n°~2013/287 le fait transmettre au CONSUPE
(Chapitre~\ref{chap:institutions}). Rédigez vos constats SI comme s'ils
devaient être repris par un vérificateur.
\end{boxmetier}

\separateur

\begin{boxresume}
\begin{itemize}
    \item \textbf{Contrôler n'est pas auditer.} Le contrôle est exercé par
          l'entité~; l'audit est un examen indépendant au profit d'un
          tiers.
    \item \textbf{Cinq éléments.} Critère, preuve, constatation,
          conclusion, opinion~; le critère est écrit avant le test.
    \item \textbf{Une assurance, pas une certitude.} Raisonnable ou limitée,
          choisie au mandat, et formulée en conséquence.
    \item \textbf{Trois types d'audit, un même SI.} L'audit des SI sert
          l'audit financier, de conformité et de performance.
    \item \textbf{Une compétence écrite.} Le décret n°~2013/287 compte le
          contrôle des SI parmi les cinq contrôles du CONSUPE.
    \item \textbf{Les jointures d'abord.} Les données de référence et les
          transmissions entre applications concentrent le risque.
    \item \textbf{La preuve change de forme.} Population dans le système,
          intégrité à établir, trois réponses possibles.
\end{itemize}
\end{boxresume}

\begin{boxexo}[title={\ding{45}\quad Exercice~1.1. Qui fait quoi~? \niveaubase}]
Pour chacune des situations suivantes, dites s'il s'agit d'un contrôle,
d'un audit, d'une inspection ou d'une vérification, et justifiez.
\begin{enumerate}[(a)]
    \item Le chef du service de la liquidation vérifie chaque matin la
          liste des mandats saisis la veille.
    \item Un vérificateur du CONSUPE examine l'exécution du budget 2025 de
          la DAF du MINÉQUIR.
    \item L'inspecteur général du ministère se rend à la DAF après un
          article de presse.
    \item SIDÉPENSE refuse un engagement supérieur au crédit disponible.
\end{enumerate}
\end{boxexo}

\begin{boxexo}[title={\ding{45}\quad Exercice~1.2. Les cinq éléments \niveaubase}]
Dans le dossier CAPNO, le fichier de paie de décembre compte 4\,312
agents, dont 12 partagent le même compte bancaire avec au moins un autre
agent. Identifiez, dans cette phrase et dans la suite que vous imaginerez,
un critère, une preuve, une constatation et une conclusion.
\end{boxexo}

\begin{boxexo}[title={\ding{45}\quad Exercice~1.3. Choisir l'assurance \niveaumoyen}]
On vous confie trois jours pour examiner la paie de la CAPNO.
\begin{enumerate}[(a)]
    \item Quel niveau d'assurance pouvez-vous donner~? Pourquoi~?
    \item Rédigez la phrase de conclusion correspondante, dans les deux
          cas~: vous n'avez rien trouvé~; vous avez trouvé 37 agents qui
          ne figurent pas au fichier du personnel.
\end{enumerate}
\end{boxexo}

\begin{boxexo}[title={\ding{45}\quad Exercice~1.4. Cartographier \niveauexpert}]
Pour le projet e-RECETTES de Bafoussam (dématérialisation des recettes),
dressez le tableau des fonctions du système, des jointures avec les
autres systèmes (trésorerie, comptabilité), et du risque principal à
chaque jointure. Dites quelles trois jointures vous testeriez en premier,
et pourquoi.
\end{boxexo}

\begin{boxinfo}
\textbf{Transition.} Ce chapitre a posé le vocabulaire et désigné
l'objet~: les systèmes qui portent les finances publiques. Le
Chapitre~\ref{chap:sci} présente ce que l'auditeur y cherche avant
tout~: le système de contrôle interne.
\end{boxinfo}
